\chapter{Методы хранения}

Необходимым условием для работы с временными рядами является организация хранения отслеживаемых данных.

В силу особенностей временных рядов и их обработки можно обозначить ключевые требования к средствам их хранения:
\begin{itemize}
    \item приоритетность скорости записи данных перед скоростью их обработки, так как операции записи преобладают над операциями чтения~\cite{poliref1};
    \item упорядоченность хранимых данных по времени, так как данные записываются, считываются и обновляются не случайным образом, а во временном порядке~\cite{poliref1};
    \item расширяемость и вместительность хранилища данных;
    \item гарантия сохранности данных. 
\end{itemize}

В данной работе в качестве методов хранения выделены и рассмотрены следующие подходы:

\begin{itemize}
    \item использование баз данных, специализированных для временных рядов -- OpenTSDB;
    \item использование нереляционных БД -- Cassandra;
    \item использование файлов. 
\end{itemize}

OpenTSDB --- технология специализированная для временных рядов, поэтому стоит уделить ей особое внимание,
так как это позволит выявить основные черты специализированных для временных рядов БД.

Нереляционная база данных является более универсальной технологией.
В качестве представителя данного класса, в данной работе рассмотрена Cassandra, так как она наиболее востребована среди подобных.

По своей сути, база данных --- это файл, обернутый в специальное программное обеспечение,
которое делает его максимально безопасным в использовании.
Иногда имеет смысл отказаться от оболочки и работать с данными напрямую.
В данном случае, рассматривается метод с использованием файлов.

\section{Хранение временных рядов в OpenTSDB}

При работе с динамическими временными рядами основной задачей является эффективная организация данных.
Поскольку временные ряды обновляются и пополняются новыми значениями, требуются оптимизированные структуры данных для быстрого доступа к информации и эффективного выполнения запросов.
Кроме того, база данных должна обеспечивать быструю и эффективную запись новых данных.

OpenTSDB (Open Time Series Database) обеспечивает хранение и агрегацию временных рядов для обеспечения масштабируемости и производительности при работе с большими объемами данных.
В OpenTSDB элемент временного ряда состоит из вещественного значения, уникального идентификатора временного ряда, метки времени и набора тегов.
Теги представляют собой символьные строки для хранения метаданных.

OpenTSDB предоставляет возможности для индексирования данных и агрегации по времени и меткам.
Это позволяет эффективно выполнять запросы к большим наборам данных временных рядов, выполнять аналитику и получать агрегированные результаты.
Данная база данных разработана с учетом горизонтального масштабирования, что позволяет легко добавлять новые узлы для увеличения общей емкости хранения и производительности.

В OpenTSDB запросы к базе данных формируются с использованием языка JSON~\cite{aliref10}.
Эти запросы поддерживают различные операции, такие как арифметические и логические выражения, фильтры, группировка, статистические и аналитические функции.
Одной из таких функций является линейная интерполяция.
Если недостаточно реальных данных для интерполяции, то результатом может быть значение <<NaN>>.
Также в OpenTSDB существуют запросы по временному диапазону, фильтрация по меткам, позволяющая выбирать временные ряды с определенными значениями метаданных, агрегация, с помощью которой есть возможность вычислить среднее значение, сумму или других статистических показателей для временных рядов.

OpenTSDB не предоставляет инструменты для интеллектуального анализа временных рядов, за исключением линейной интерполяции, но допускает расширения от сторонних разработчиков, добавляющие указанную функциональность (например, библиотека R2Time~\cite{aliref11}, реализованная на языке программирования~R).

\section{Хранение временных рядов в Cassandra}

Apache Cassandra --- базы данных, которая является:

\begin{itemize}
    \item \textit{распределенной} --- обеспечивает работу на нескольких машинах, в дальнейшем называемых узлами, которые пользователем представляются как единое целое~\cite{myref3};
    \item \textit{децентрализованной} --- все узлы являются равноправными~\cite{myref3};
    \item \textit{эластично масштабируемой} --- обеспечивает возможность добавлять и убирать узлы, не прерывая работы системы.
Данное свойство является специальным для горизонтальной масштабируемости --- механического увеличения пропускной способности оборудования путем добавления памяти в уже имеющуюся машину~\cite{myref3};
    \item \textit{отказоустойчивой} --- предполагает замену поврежденных узлов без простоя, за счет их репликации на несколько узлов~\cite{myref3}
(репликация представляет собой набор технологий копирования и распространения данных и объектов между узлами, а также включает синхронизацию этих узлов для поддержания их согласованности
(согласованность означает, что операция чтения всегда возвращает последнее записанное значение)~\cite{myref3, myref4});
    \item \textit{допускающей настройку согласованности} --- существует возможность настроить количество узлов, которые должны подтвердить завершение операции обновления данных~\cite{myref3}.
\end{itemize}
Основными структурами данных в Cassandra являются:
\begin{itemize}
\item \textit{столбец} --- пара имя-значение;
\item \textit{строка} --- контейнер для упорядоченной коллекции столбцов на которые можно сослаться по первичному ключу;
\item \textit{таблица} --- контейнер для упорядоченной коллекции строк;
\item \textit{пространство ключей} --- контейнер таблиц;
\item \textit{кластер} --- контейнер пространств ключей, расположенных в одном или нескольких узлах~\cite{myref3}.
\end{itemize}

Таким образом, столбец --- самая элементарная структурная единица в модели данных Cassandra, и ему присущи следующие характеристики:
имя; значение; тип; временная метка, генерируемая при каждом обновлении столбца; время жизни, указывающее срок хранения значения столбца~\cite{myref3}.

Уточнив основные понятия базы данных, можно рассмотреть особенности Cassandra применительно к хранению временных рядов.

Модель данных организует хранение временного ряда, где
\textit{ключ для строки} --- имя для нашего измерения,
\textit{имя столбца} --- время измерения,
\textit{значение столбца} --- измеренное значение~\cite{myref5}.
Данная структура позволяет выбрать как индивидуальное значение, так и набор значений для заданного интервала какого-либо измерения~\cite{myref5}.

Кроме того, для каждого столбца хранится время жизни, позволяющее автоматизировать процесс удаления устаревших данных, а также временная метка, которая используется с целью:

\begin{itemize}
\item разрешения конфликтов при изменении одного и того же значения со стороны нескольких клиентов~\cite{myref3},
\item организации автоматической сортировки данных по времени~\cite{myref5}. 
\end{itemize}

В случае обработки временных рядов скорость поступления данных имеет большее значение, нежели скорость их изменения,
потому что информация однородна и фактически обрабатывается только небольшая часть из всего объема хранимой информации~\cite{myref5}.
Данные особенности позволяют выделить Cassandra, как нереляционную БД (скорость добавления данных в NoSQL системах выше, чем в SQL системах)~\cite{myref5}.
Cassandra является распределенной и децентрализованной, вследствие этого можно выделить следующие ее аспекты:
\begin{itemize}
\item отказоустойчивость --- гарантия сохранности информации за счет репликации~\cite{myref6};
\item масштабируемость --- увеличение объема хранимых данных возможно посредством добавления узла в кластер~\cite{myref6};
\item слабая согласованность данных --- отсутствие контроля актуальности читаемой информации~\cite{myref3}.
Данная особенность может быть скорректирована настройкой согласованности. 
\end{itemize}

\section{Хранение временных рядов в  файлах}

Временные ряды могут храниться в текстовых, CSV и JSON файлах.
В данном случае все взаимодействия с данными, а именно распределение памяти для хранения данных, открытие,
чтение и закрытие файлов, их защита и так далее, выполняются специалистом по анализу данных «вручную».

\textit{Преимущества} данного способа, по сравнению с использованием баз данных, следующие~\cite{poliref1}:

\begin{itemize}
\item не требует значительных временных затрат на ввод/вывод данных, так как данная операция производится посредством чтения/изменения файла;
\item существует возможность максимально увеличить сжатие данных, что позволяет сократить требуемый объем памяти для хранения временного ряда;
\item обеспечивает хранение различных типов временных рядов, включая данные с разной частотой и разрешением;
\item используется для обмена данными между различными системами;
\item применяется для резервного копирования данных.
\end{itemize}

\textit{Недостатками} данного способа, по сравнению с использованием баз данных, являются~\cite{poliref1}:

\begin{itemize}
\item снижение производительности при увеличении объема хранимых данных;
\item отсутствие возможности масштабировать систему хранения;
\item уязвимость системы к несанкционированным доступам и изменениям;
\item отсутствие эффективной организации работы с данными при увеличении количества пользователей;
\item необходимость обеспечивать и поддерживать структуру и целостность данных.
\end{itemize}

\section{Сравнение методов хранения временных рядов}

Результаты классификации методов хранения временных рядов приведены в таблице~\ref{tab_mem}.
\begin{table}[h!]
    \caption{Классификация методов хранения временных рядов}
    \label{tab_mem}
    \centering
    \begin{tabular}{||p{4.2cm}||p{3cm}|p{5cm}|p{3cm}||}
        \hline
        & \multicolumn{3}{c|}{Система хранения} \\ \cline{2-4}
        Параметр & Cassandra & OpenTSDB & файлы\\
        \hline \hline
        {автоматическое удаление устаревших данных} & есть & нет & нет \\
        \hline
        {автоматическая сортировка данных по времени} & есть & нет & нет \\
        \hline
        {сравнение скорости поступления и обработки данных} &  {скорость поступления данных выше} &  {скорость поступления данных выше~\cite{itiref2}} &  {скорость поступления и обработки данных не имеет значительных различий}\\[8ex]
        \hline
        {отказоустойчивость} & есть & есть & нет\\
        \hline
        {масштабируемость} &{гибкая масштабируемость} &  {горизонтальная масштабируемость~\cite{itiref1}} & нет\\
        \hline
        {инструменты для интеллектуального анализа временных рядов} & нет &  {линейная интерполяция, допускаются расширения от сторонних разработчиков} & нет\\
        \hline
        {возможность сократить объем требуемой памяти} & нет & нет & есть\\
        \hline
        {встроенный модуль для визуализации данных} & нет &  {есть встроенный инструмент построения линейных графов на основе нескольких метрик и меток~\cite{itiref1}} & нет\\
        \hline
        {составляющие единицы данных} &  {временная метка, значение, время жизни} &  {любое количество значений, временная метка} & {значение, отметку времени}\\
        \hline \hline
    \end{tabular}
\end{table}
\clearpage

В силу того, что OpenTSDB является надстройкой над некоторой базой данных, в том числе и Cassandra~\cite{itiref3},
становится очевидным его преимущество перед ней в возможности автоматизированного интеллектуального анализа временных рядов,
визуализации данных и т. п..
Однако, автоматизация сортировки данных по времени и удаления устаревшей информации характерна только для Cassandra.
С другой стороны, в OpenTSDB возможно хранить несколько значений для одной временной метки.

Несмотря на большое количество недостатков файловой организации по сравнению с базами данных,
стоит отметить некоторые их особые области применения: обмен данными между различными системами и резервное копирование.
Возможность максимально сжать информацию делает предложенный метод, применительно к конкретным задачам,
конкурентным по отношению к базам данных, при этом сохраняется перспектива использования файлов совместно с БД.

Таким образом, в работе были рассмотрены три различных подхода к хранению временных рядов,
выделены особенности их использования, на примере одного из представителей каждого вида и
проведена классификация методов по параметрам, определенных требованиями к хранению временных рядов.

